% examples/industry/resume-english.tex
%
% The supported v0.1.0 English industry résumé example. It is the reference
% document for careerdossier-resume.cls: a single-column, monochrome, US Letter
% résumé built entirely from the public API (\CDossierSetup, \MakeCDossierHeader,
% \CDossierSection, \CDossierSubsection, CDossierEntry, CDossierItemize).
%
% It is also the example that demonstrates `contact-labels' (see below). It is
% deliberately left on the default *untagged* path, because that is the output
% the label is there to fix; the tagged path is demonstrated instead by
% examples/industry/letter-industry.tex, which shares this profile.
%
% Compile with LuaLaTeX:
%   latexmk -lualatex -interaction=nonstopmode -halt-on-error \
%     examples/industry/resume-english.tex
% or:
%   lualatex examples/industry/resume-english.tex
%
% Run from the repository root so the class and shared packages resolve, or add
% the repository root to TEXINPUTS.

\documentclass[fontsize=11pt, margin=narrow]{careerdossier-resume}

% examples/profiles/profile-english.tex
%
% Shared English profile metadata for the CareerDossierTeX examples. It carries
% personal contact data only; document-specific content lives in the resume and
% letter files that input this profile. Keeping the data here means the same
% contact information is not copied across every example.
%
% This file is meant to be \input by an example document; it is not a standalone
% compilable file.

\CDossierSetup{
  name     = {Ada Lovelace},
  headline = {Data Scientist},
  email    = {ada.lovelace@example.com},
  phone    = {+1 416 555 0142},
  location = {Toronto, Ontario, Canada},
  website  = {example.com},
  %% linkedin, github, and scholar accept either form, and the two below are
  %% deliberately written differently to show both. A bare handle is expanded
  %% to the service's canonical address, so `ada-lovelace' and
  %% `linkedin.com/in/ada-lovelace' produce the same page and the same link;
  %% a value that already carries a `/' or a `.' is used exactly as written.
  %% See docs/API.md, "Accepted forms for the web-profile keys", for the forms
  %% that build but link nowhere.
  linkedin = {ada-lovelace},
  github   = {github.com/ada-lovelace},
}


% `contact-labels' prefixes a short identifying label -- Email, Phone,
% Website -- to the three contact fields whose value does not say what it is.
% It is an option key rather than profile data, so it belongs here in the
% document and not in the shared profile file.
%
% Why a résumé would want it: an e-mail address and a website are announced as
% links by a screen reader, but a phone number is announced as bare digits with
% nothing conveying what it is, and plain-text extraction and ATS parsers see
% no link annotations at all. A visible label is the one mechanism that reaches
% every consumer, including the untagged output this example produces.
%
% It is opt-in and off by default; deleting this line restores the unlabelled
% contact line. See docs/API.md, "Contact-field labels (contact-labels)".
\CDossierSetup{ contact-labels = true }

\begin{document}

\MakeCDossierHeader

\CDossierSection{Summary}

Data scientist with experience designing reliable analytical platforms,
automated reporting pipelines, and accessible public-facing services. Skilled in
translating operational requirements into maintainable systems, measurable
service standards, and clear technical documentation.

\CDossierSection{Experience}

% \CDossierSubsection is the level between a section and an entry. It groups
% entries that belong together without promoting each group to a ruled section
% of its own, which would flatten the hierarchy the grouping is meant to show.
% A section that opens directly with a subsection, as this one does, leaves the
% same gap as one opening directly with an entry.
\CDossierSubsection{Public Sector}

\begin{CDossierEntry}[
  title        = {Senior Data Scientist},
  organization = {Government Digital Services Laboratory},
  location     = {Ottawa, ON / Remote},
  dates        = {2024--Present}
]
  \begin{CDossierItemize}
    \item Lead the delivery of reporting services supporting high-volume public
          programs and internal operational dashboards.
    \item Consolidated several ad hoc analyses into a shared, governed pipeline,
          reducing duplicated maintenance across program teams.
    \item Established automated validation for data quality, reproducibility, and
          archival output.
  \end{CDossierItemize}
\end{CDossierEntry}

\CDossierSubsection{Research and Nonprofit}

\begin{CDossierEntry}[
  title        = {Data Scientist},
  organization = {Civic Technology Research Group},
  location     = {Toronto, ON},
  dates        = {2021--2024}
]
  \begin{CDossierItemize}
    \item Built models and processing pipelines that improved reliability during
          seasonal demand peaks and large reporting cycles.
    \item Introduced contract tests and operational dashboards that shortened
          incident investigation.
  \end{CDossierItemize}
\end{CDossierEntry}

\CDossierSection{Selected Projects}

\begin{CDossierEntry}[
  title = {Operational Reporting Modernization},
  dates = {2023--2024}
]
  \begin{CDossierItemize}
    \item Replaced spreadsheet-driven monthly reporting with a reproducible
          publication process and reconciled metric definitions.
  \end{CDossierItemize}
\end{CDossierEntry}

\CDossierSection{Education}

\begin{CDossierEntry}[
  title        = {M.Sc. in Software Engineering},
  organization = {McGill University},
  location     = {Montréal, QC},
  dates        = {2019--2021}
]
\end{CDossierEntry}

\begin{CDossierEntry}[
  title        = {B.A.Sc. in Computer Engineering, Honours},
  organization = {University of Waterloo},
  location     = {Waterloo, ON},
  dates        = {2014--2019}
]
\end{CDossierEntry}

\CDossierSection{Technical Skills}

\textbf{Programming:} Python, SQL, R, Bash \\
\textbf{Platforms:} Linux, Docker, PostgreSQL, cloud CI/CD \\
\textbf{Practices:} Reproducible analysis, automated testing, data validation,
observability

\end{document}
